\documentclass[11pt,a4paper]{letter}
\usepackage[utf8]{inputenc}
\usepackage[T1]{fontenc}
\usepackage[french]{babel}
\usepackage{lmodern}
\usepackage{geometry}
\geometry{left=2cm,right=2cm,top=2cm,bottom=2cm}
\usepackage{hyperref}
\usepackage{xcolor}

\definecolor{titleblue}{RGB}{0,51,140}
\hypersetup{
    colorlinks=true,
    urlcolor=titleblue,
    linkcolor=titleblue
}

\begin{document}

\begin{flushleft}
\textbf{Ismail AISSAOUI}\\
Chlef, Algérie\\
Tél : +213 660 70 77 96\\
Email : \href{mailto:ismail.aissaoui.pro@gmail.com}{ismail.aissaoui.pro@gmail.com}\\
LinkedIn : \href{https://linkedin.com/in/aissaoui-ismail-6a77a92a7}{linkedin.com/in/aissaoui-ismail-6a77a92a7}
\end{flushleft}

\vspace{1cm}

\begin{flushright}
\textbf{À l'attention de Mounira Harzallah}\\
\textbf{Carito Guziolowski et André Ndobo}\\
Laboratoires LS2N et LPPL, Nantes\\
\end{flushright}

\vspace{1cm}

\textbf{Objet :} Candidature à la thèse « Détection explicable des biais discriminatoires par construction d'ontologies et LLM augmentés par des connaissances »

\vspace{0.5cm}

Madame, Monsieur,

Actuellement étudiant en dernière année d'ingénierie et de Master en Intelligence Artificielle et Data Science à l'École Nationale Supérieure d'Informatique (ESI-SBA), c'est avec un vif intérêt que je vous soumets ma candidature pour l'offre de thèse portant sur la modélisation formelle et la détection explicable des biais discriminatoires, sous la direction de Mounira Harzallah, Carito Guziolowski et André Ndobo.

Le croisement entre l'intelligence artificielle responsable, l'extraction de connaissances, et la psychologie sociale m'interpelle particulièrement. L'utilisation des grands modèles de langage (LLM) et du Web sémantique (RAG, ontologies, graphes de connaissances) pour lutter contre les biais cognitifs et les discriminations dans les processus de décision répond à un enjeu sociétal critique. La perspective de contribuer à l'élaboration d'un cadre computationnel formel capable de rendre les décisions IA plus équitables, transparentes et explicables s'inscrit parfaitement dans ma démarche scientifique et mes aspirations professionnelles.

Au cours de ma formation et de mes expériences de recherche, j'ai développé de solides compétences en Deep Learning, en traitement du langage naturel (NLP), ainsi que dans l'architecture et l'optimisation des réseaux de neurones. J'ai notamment conçu un système de recommandation de recrutement exploitant les réseaux de neurones sous forme de graphes (GNN) et les embeddings sémantiques (SBERT), afin de modéliser avec précision les relations complexes entre offres d'emploi et profils de candidats. Ces travaux sur la représentation sémantique et la modélisation relationnelle rejoignent directement les enjeux de construction automatique d'ontologies et d'identification des relations entre biais que propose votre projet de thèse.

Par ailleurs, mon expérience récente en tant qu'ingénieur de recherche chez Sonatrach, où j'ai architecturé des Jumeaux Numériques intégrant des contraintes complexes de physique formelle, atteste de ma capacité à lier les données d'apprentissage profond à des systèmes de connaissances structurées a priori. Maîtrisant Python, PyTorch et les architectures avancées (Transformers, LLM, architectures séquentielles), je suis techniquement armé pour explorer l'intégration des LLM et des méthodes RAG à la structuration ontologique requise par le projet. 

Je suis profondément motivé par la perspective d'évoluer au sein de l'environnement interdisciplinaire des laboratoires LS2N et LPPL, alliant l'informatique à la psychologie sociale, afin de concevoir un système qui sensibilise tout en fournissant une justification explicable et modérée. Mon autonomie, ma rigueur scientifique et mon grand intérêt pour l'éthique de l'IA seront, je l'espère, de solides atouts pour la Chaire « Origines et mécanismes des discriminations ».

Dans l'attente d'un éventuel entretien qui me permettrait de vous détailler mon parcours et mon enthousiasme pour ce projet, je vous prie d'agréer, Madame, Monsieur, l'expression de mes salutations distinguées.

\vspace{1cm}
\textbf{Ismail AISSAOUI}

\end{document}
